\subsection{Twisting and distance bounds}
\label{sec:twisted-flow-distance-dependence}

A twisted Chern--Ricci flow is the equation $\partial_t\omega=-\operatorname{Ric}(\omega)+\zeta$, where the twist $\zeta$ is a fixed smooth closed real $(1,1)$-form. Diameter bounds along such flows cannot depend only on the maximal time and on the initial Monge--Amp\`ere density: the twist and the length of the time interval enter the geometry, even for flat K\"ahler metrics.

\Needspace{14\baselineskip}
\begin{example}\label{ex:twisted-flow-distance-dependence}\index{Chern--Ricci flow!dependence on the twist}
Let $E=\C/(\mathbb Z+\mathrm i\mathbb Z)$, let $X=E^n$ with $n\geq2$, and let $\eta_j=(\mathrm i/2)\mathrm dz_j\wedge\mathrm d\bar z_j$ be the flat forms pulled back from the factors. Put $\omega_0=\sum_{j=1}^n\eta_j$. All the flows below start at $\omega_0$, with potential zero and initial density one relative to $\omega_0^n$. In particular, hypotheses on the initial data expressed through a decomposition $\mathrm e^{\psi^+-\psi^-}$ of the density, with quasi-plurisubharmonic functions $\psi^\pm$ subject to fixed $L^p$ bounds, hold uniformly for all these flows, with $\psi^+=\psi^-=0$.

For each $M>0$, the twist $\zeta_M=M\eta_1-\eta_2$ gives the explicit solution
\[
 \omega_M(t)=(1+Mt)\eta_1+(1-t)\eta_2+
                         \sum_{j=3}^n\eta_j,
 \qquad 0\leq t<1,
\]
of $\partial_t\omega=-\operatorname{Ric}(\omega)+\zeta_M$. Its maximal smooth existence time is exactly one, while
\[
 \operatorname{diam}(X,\omega_M(1/2))
 \geq \sqrt{1+M/2}\,\operatorname{diam}(E,\eta_1).
\]
For the single fixed twist $\zeta=\omega_0$, the solution is $(1+t)\omega_0$ for all $t\geq0$, and its diameter grows as $\sqrt{1+t}\,\operatorname{diam}(X,\omega_0)$.
\end{example}
\begin{proof}
Each of these metrics has constant coefficients, so its Chern--Ricci form vanishes; differentiation in $t$ proves the flow equations, and positivity of the coefficients gives the stated existence intervals. In the first family, pair a hypothetical positive representative at a time $t\geq1$ with the closed positive form $\eta_1\wedge\eta_3\wedge\cdots\wedge\eta_n$, where the empty product is $1$ when $n=2$. The potential term integrates to zero, and the pairing is a positive fixed multiple of $1-t$, so it cannot be strictly positive. Integrating the flow equation shows that every smooth solution differs from the affine background by $\ddc$ of a smooth function, so no smooth positive continuation beyond $t=1$ is possible.

Projection to the first elliptic factor does not increase lengths once the metric of that factor is scaled by $1+Mt$. Taking two points that differ only in the first factor proves the lower bound for the diameter. For the second flow, every path length is multiplied by $\sqrt{1+t}$, which proves the exact formula for the diameter.
\end{proof}

On a fixed finite time interval, a common Dini modulus for the potentials and a uniform bound for the affine backgrounds do give control of distances and diameters, by \cref{cor:dini-family-distance-limits}. The example shows why such bounds must be imposed when the twist varies or the time interval grows.
